\section*{The Binomial Theorem}
\addcontentsline{toc}{section}{The Binomial Theorem}

\begin{theorem}[Binomial Theorem]
For any integer \(n \ge 0\) and real numbers \(x,y\),
\[
(x+y)^n
=
\sum_{k=0}^{n} \binom{n}{k} x^{\,n-k} y^{\,k},
\]
where the coefficients
\[
\binom{n}{k}
=
\frac{n!}{k!(n-k)!}
\]
are called \emph{binomial coefficients}.
\end{theorem}

\begin{remark}
The binomial theorem describes how powers of a sum expand into a sum of
products. The coefficients count the number of ways to choose \(k\) terms
of \(y\) from the \(n\) factors in the product
\[
(x+y)(x+y)\cdots(x+y).
\]
Thus the binomial coefficients arise naturally from combinatorics as the
number of ways to select \(k\) objects from a set of \(n\) objects.
\end{remark}

\subsection*{Basic Properties of Binomial Coefficients}

\begin{align*}
\binom{n}{0} &= 1, \\
\binom{n}{n} &= 1, \\
\binom{n}{k} &= \binom{n}{n-k}.
\end{align*}

They also satisfy the recursive identity
\[
\binom{n+1}{k}
=
\binom{n}{k}+\binom{n}{k-1},
\]
which produces Pascal's triangle.

\begin{center}
\[
\begin{array}{ccccccccccC}
&&&&&1\\[6pt]
&&&&1&&1\\[6pt]
&&&1&&2&&1\\[6pt]
&&1&&3&&3&&1\\[6pt]
&1&&4&&6&&4&&1\\[6pt]
1&&5&&10&&10&&5&&1
\end{array}
\]
\end{center}



\begin{center}
\begin{tikzpicture}[scale=1]
\foreach \n in {0,...,6}
{
  \foreach \k in {0,...,\n}
  {
    \node at (\k-\n/2,-\n) {$\binom{\n}{\k}$};
  }
}
\end{tikzpicture}
\end{center}

\begin{remark}
Pascal's triangle encodes the binomial coefficients
\[
\binom{n}{k}.
\]
Each entry is obtained from the sum of the two entries above it:
\[
\binom{n}{k}=\binom{n-1}{k-1}+\binom{n-1}{k}.
\]
These coefficients appear throughout discrete calculus, especially
in formulas for higher finite differences and in Newton's forward
difference expansion.
\end{remark}


\subsection*{Connection with Discrete Calculus}

The binomial theorem plays a fundamental role in discrete calculus because
difference operators naturally produce binomial coefficients.

\begin{proposition}
For any function \(f\),
\[
\Delta^n f(x)
=
\sum_{k=0}^{n} (-1)^{\,n-k}
\binom{n}{k}
f(x+k).
\]
\end{proposition}

\begin{remark}
This identity shows that higher finite differences are governed by the
same combinatorial structure that appears in the binomial theorem.
In discrete calculus, binomial coefficients arise naturally in
\begin{itemize}
\item formulas for higher differences,
\item Newton's forward difference expansion,
\item polynomial representations in the falling factorial basis,
\item discrete Taylor series.
\end{itemize}

Because of this, the binomial theorem serves as one of the central
algebraic tools underlying the theory of finite differences.
\end{remark}

\subsection*{Example}

Applying the difference operator repeatedly to a function yields

\[
\begin{aligned}
\Delta f(x) &= f(x+1)-f(x), \\
\Delta^2 f(x) &= f(x+2)-2f(x+1)+f(x), \\
\Delta^3 f(x) &= f(x+3)-3f(x+2)+3f(x+1)-f(x),
\end{aligned}
\]

which mirrors the binomial expansion of

\[
(1-1)^n.
\]

Thus the binomial coefficients determine the pattern of coefficients
appearing in higher-order finite differences.